\subsection*{A concrete tensor}

Let $e_1=(1,0)^{\mathsf T}$ and $e_2=(0,1)^{\mathsf T}$. Start with
\[
 D_{(2,1)}=e_1\otimes e_1\otimes e_2
          +e_1\otimes e_2\otimes e_1
          +e_2\otimes e_1\otimes e_1.
\]
This is a $2\times2\times2$ array: its entries at $(1,1,2)$, $(1,2,1)$
and $(2,1,1)$ are one, and all other entries are zero. It sums all
arrangements of two occurrences of the first symbol and one of the
second. A general Dicke tensor does the same for an arbitrary positive
multiplicity tuple $\mathbf d$.

\begin{center}
\begin{tabular}{@{}lp{0.51\textwidth}l@{}}
\toprule
Symbol & Meaning & This example\\
\midrule
$n$ & Number of symbols; dimension of each original factor & $2$\\
$\mathbf d$ & Positive symbol multiplicities & $(2,1)$\\
$k=\sum_a d_a$ & Number of original tensor factors & $3$\\
$N$ & Number of copies of the whole tensor & Any positive integer\\
$q$ & Size of the extracted diagonal tensor & An integer to maximize\\
$p=\mathbf d/k$ & Symbol-frequency distribution & $(2/3,1/3)$\\
\bottomrule
\end{tabular}
\end{center}

\subsection*{From diagonal extraction to an entropy-determined rate}

The target of a restriction is the unit tensor
$\Delta_q=\sum_{j=1}^q e_j^{\otimes k}$. For example,
$\Delta_2=e_1\otimes e_1\otimes e_1+e_2\otimes e_2\otimes e_2$
when $k=3$. This is a target diagonal structure, not another expression
for $D_{(2,1)}$. Subrank asks for the largest such $q$ obtainable by
independent complex-linear maps on the factors. For two factors this
is matrix rank; with more factors it measures the size of a diagonal
that can be extracted by these separate maps.

Taking $N$ copies groups corresponding positions only. Thus two copies
of the example form a $4\times4\times4$ array, by the regrouping
\[
 (e_a\otimes e_b\otimes e_c)\otimes
 (e_{a'}\otimes e_{b'}\otimes e_{c'})
 \longmapsto
 (e_a\otimes e_{a'})\otimes(e_b\otimes e_{b'})\otimes
 (e_c\otimes e_{c'}).
\]
There remain three factors. In general there remain $k$ factors, each
of dimension $n^N$; we never merge different original modes.

For $q_N=Q(D_{\mathbf d}^{\boxtimes N})$, the theorem identifies the
leading exponential growth by the entropy of the symbol frequencies:
\[
 \lim_{N\to\infty}q_N^{1/N}=2^{\HH_2(p)},\qquad
 \log_2 q_N=N\HH_2(p)+o(N),\qquad
 \HH_2(p)=-\sum_a p_a\log_2 p_a.
\]
In the example,
\[
 \HH_2(2/3,1/3)\approx0.9183,\qquad
 2^{\HH_2(2/3,1/3)}=\frac{3}{2^{2/3}}\approx1.8899.
\]
The growth base need not be an integer. This does not say that any
finite $q_N$ equals $(1.8899)^N$, or that the single-copy subrank equals
that noninteger. The limit is taken for each fixed tuple $\mathbf d$.
The example belongs to the W-tensor family treated by
\citet[Theorem~11]{VC2015}; it illustrates the question, rather than
the extent of the all-multiplicity conclusion.

The general entropy formula was proposed as an entanglement-conversion
rate in \citet[Section~VI, equation~(40)]{VC2015} and restated for
ordinary asymptotic subrank in \citet[Conjecture~4.26]{OV2025}.
Theorem~\ref{thm:main} settles it for every positive multiplicity tuple.
It also proves that coordinate selection, followed by independent
relabelling, attains the same exponential rate as arbitrary independent
linear maps. The question settled is therefore the exact leading
exponent of extractable diagonal structure throughout this family.
The source conventions and their comparison are detailed in
Section~\ref{sec:history}.

\subsection*{Proof roadmap: symmetry controls entropy, entropy controls selection}

The central entropy statement is Lemma~\ref{lem:invariant}. Let $E$
be a positive-dimensional finite-dimensional vector space over $\Q$,
let $X$ be an $E$-valued random variable with finite support, and let a
group act $\Q$-linearly and irreducibly on $E$. Assume that $gX$ has
the same law as $X$ for every group element $g$. Then every rational
linear subspace $V\leq E$ satisfies
\[
 \HH_2(X\mid X\bmod V)
 \leq \frac{\dim V}{\dim E}\,\HH_2(X).
\]
Here $X\bmod V$ is the coset label. The base of the logarithm is
immaterial to this inequality; Section~\ref{sec:entropy} uses natural
logarithms throughout its proof.

The function $V\mapsto\HH(X\mid X\bmod V)$ is supermodular on the
subspace lattice. Among nonzero subspaces maximizing entropy per dimension,
choose one of largest dimension. Invariance and supermodularity force
it to be invariant under the group, so irreducibility makes it all of
$E$. Applied to the permutation-invariant conditional distributions
obtained by fixing one tensor coordinate, this bounds the entropy of
each collision fiber in proportion to its dimension.

This entropy bound becomes an exact bound on the number of possible
partners of a candidate of one full joint type. The partner may lie
anywhere in the original support. Rank-dependent hashing and an
average-free set then control simultaneous survival of colliding
edges. Global isolation leaves a diagonal support obtainable by local
coordinate deletion and relabelling. Isolation is checked against
\emph{all original support edges}, so no uncounted mixed edge remains.
Sections~\ref{sec:fibers}--\ref{sec:lower-bound} give the counts,
modular-rank argument, construction and asymptotic passage in full.

For the opposite inequality, a slice cover gives
$Q(D_{\mathbf d}^{\boxtimes N})\leq k\,2^{N\HH_2(p)}$ for every
$N\geq1$, allowing arbitrary independent linear maps
(Section~\ref{sec:upper}). Together, the bounds identify the entropy
rate. The existence of the limit and the two-factor endpoint are
included in Section~\ref{sec:limits}.

\subsection*{Precise statement and conventions}
